\documentclass{article}

\usepackage{booktabs}
\usepackage{tabularx}
\newcolumntype{L}{>{\raggedright\arraybackslash}X}

\title{Development Plan\\\progname}

\author{\authname}

\date{}

\input{../Comments.text}
\input{../Common.text}

\begin{document}

\maketitle

\begin{table}[hp]
\caption{Revision History} \label{TblRevisionHistory}
\begin{tabularx}{\textwidth}{llX}
\toprule
\textbf{Date} & \textbf{Developer(s)} & \textbf{Change}\\
\midrule
Date1 & Name(s) & Description of changes\\
Date2 & Name(s) & Description of changes\\
... & ... & ...\\
\bottomrule
\end{tabularx}
\end{table}

\newpage{}

\wss{Put your introductory blurb here.  Often the blurb is a brief roadmap of
what is contained in the report.}

\wss{Additional information on the development plan can be found in the
\href{https://gitlab.cas.mcmaster.ca/courses/capstone/-/blob/main/Lectures/L02b_POCAndDevPlan/POCAndDevPlan.pdf?ref_type=heads}
{lecture slides}.}

\section{Confidential Information?}

\wss{State whether your project has confidential information from industry, or
not.  If there is confidential information, point to the agreement you have in
place.}

\wss{For most teams this section will just state that there is no confidential
information to protect.}
\section{IP to Protect}

\wss{State whether there is IP to protect.  If there is, point to the agreement.
All students who are working on a project that requires an IP agreement are also
required to sign the ``Intellectual Property Guide Acknowledgement.''}

\section{Copyright License}

\wss{What copyright license is your team adopting.  Point to the license in your
repo.}

\section{Team Meeting Plan}

\indent In-person team meetings will be held weekly on Thursdays from 3:30pm to 4:20pm. The purpose of these meetings will be to discuss upcoming deliverables, progress updates, and potential blockers. They will also be used to consolidate any questions that team members may have for the instructional team and project supervisors.

\indent There will also be biweekly meetings held with the project supervisors on Fridays from 9:00am to 10:00am. They are planned to be in-person but may be shifted to virtual Microsoft Teams meetings based on availability. These will be used to ensure the project is in line with the advisors’ requirements, and to address any questions that the development team might have.

\indent The meeting chair (Sarah Dorfman) will prepare an agenda for the main topics of discussion and distribute it to the team at least one day prior to any team meetings to give all members time to prepare. All members are encouraged to take notes during the meeting. After the meeting, the team will collectively assign action items to each other. If there is a supervisor meeting planned for that week, the team will also plan a corresponding agenda. 


\section{Team Communication Plan}

\wss{Issues on GitHub should be part of your communication plan.}

\section{Team Member Roles}

\wss{You should identify the types of roles you anticipate, like notetaker,
leader, meeting chair, reviewer.  Assigning specific people to those roles is
not necessary at this stage.  In a student team the role of the individuals will
likely change throughout the year.}

\section{Workflow Plan}

\begin{itemize}
	\item How will you be using git, including branches, pull request, etc.?
	\item How will you be managing issues, including template issues, issue
	classification, etc.?
  \item Use of CI/CD
\end{itemize}

\section{Project Decomposition and Scheduling}

\begin{itemize}
  \item How will you be using GitHub projects?
  \item Include a link to your GitHub project
\end{itemize}

\wss{How will the project be scheduled?  This is the big picture schedule, not
details. You will need to reproduce information that is in the course outline
for deadlines.}

\section{Proof of Concept Demonstration Plan}

At a high level, this application will listen to child speech, process it into its phonetic representation, and analyze this to give feedback to SLPs. The greatest risks associated with this lie in the intermediate audio processing stage:

\begin{enumerate}
  \item \textbf{Converting Child Speech To Phonetic Representation -} Since children are in a developmental phase, their speech patterns vary significantly when compared to adults. Therefore, it may be challenging to ensure an acceptable accuracy rate when converting child speech into phonetic representations.
  \item \textbf{Ensuring An Acceptable Processing Time -} For this application to be viable in a clinical setting, it will need to give real time feedback to the SLP during an appointment on their computers. As per the project supervisors, existing child speech interpretation models are quite slow, so it may be difficult for this one to process the data fast enough to give the clinician feedback in less time than it would have taken for them to reach a deduction themself.
\end{enumerate}

\indent To ensure these risks can be overcome, a suitable training dataset must be found online. To ensure speech to phonetic conversion accuracy, there must be sufficient data to properly train the model. It is important to note that this data must be publicly available to comply with PHIPA patient confidentiality guidelines. 

\indent Once this data is sourced, we will need to develop the actual model that takes child speech and converts it into a readable format with the corresponding phoneme representations. If this model meets the established accuracy benchmarks, it will be feasible to continue designing this application for pediatric speech purposes. Moreover, it is fair to assume that speech therapy offices will use typical office computers that are not equipped with powerful hardware. Therefore, it is necessary to show that the model can meet the benchmark processing time on a device of similar specifications.


\section{Expected Technology}

\wss{What programming language or languages do you expect to use?  What external
libraries?  What frameworks?  What technologies.  Are there major components of
the implementation that you expect you will implement, despite the existence of
libraries that provide the required functionality.  For projects with machine
learning, will you use pre-trained models, or be training your own model?  }

\wss{The implementation decisions can, and likely will, change over the course
of the project.  The initial documentation should be written in an abstract way;
it should be agnostic of the implementation choices, unless the implementation
choices are project constraints.  However, recording our initial thoughts on
implementation helps understand the challenge level and feasibility of a
project.  It may also help with early identification of areas where project
members will need to augment their training.}

Topics to discuss include the following:

\begin{itemize}
\item Specific programming language
\item Specific libraries
\item Pre-trained models
\item Specific linter tool (if appropriate)
\item Specific unit testing framework
\item Investigation of code coverage measuring tools
\item Specific plans for Continuous Integration (CI), or an explanation that CI
  is not being done
\item Specific performance measuring tools (like Valgrind), if
  appropriate
\item Tools you will likely be using?
\end{itemize}

\wss{git, GitHub and GitHub projects should be part of your technology.}

\section{Use of AI and LLMs}

\wss{This section is a continuation of the previous section.  The role of LLMs
in your project as important enough to warrant their own section.  
How will your team be using LLMs for code and documentation development?  
What tools will you be using?  How will you verify the output of your LLMs?}

\section{Coding Standard}

\wss{What coding standard will you adopt?}

\newpage{}

\section{Appendix --- Generative AI Use Disclosure}

All writing was done by the team members by hand. All research and facts were obtained through academic sources, or by consulting the project supervisors unless explicitly stated below.

\subsection{AI Tools Used}

\begin{table}[hp]
\caption{List of AI Tools Used} \label{TblRevisionHistory}
\begin{tabularx}{\textwidth}{llX}
\toprule
\textbf{Tool Number} & \textbf{Name} & \textbf{Model} \\
\midrule
1 & Claude & Sonnet 5 \\
2 & ChatGPT & GPT 5 Thinking Mini \\
\bottomrule
\end{tabularx}
\end{table}

\subsection{How and Where Used}

\begin{table}[hp]
\caption{Usage Details for All AI Tools Used} \label{TblRevisionHistory}
\begin{tabularx}{\textwidth}{
  >{\hsize=0.5\hsize}L
  >{\hsize=0.5\hsize}L
  >{\hsize=1\hsize}L
}
\toprule
\textbf{Tool Number} & \textbf{Location} & \textbf{Details} \\
\midrule
1 & Bibliography & Used to insert a reference to the references.bib file. \\
2 & Problem Statement and Goals; (1.3) Stakeholders & Instructed the AI to think from the perspective of a patient’s parent when assessing stakeholders: highlighted privacy concerns as another consideration for this stakeholder. \\
\bottomrule
\end{tabularx}
\end{table}

\subsection{Verification of AI-Generated Content}

As none of the AI usage was for factual content, verification was very simple. For tool 1, we simply verified the resulting citation for correctness in the final output. For tool 2, it was used as more of a brainstorming utility, so only unconsidered elements that made logical sense were included. 

\newpage{}

\section*{Appendix --- Reflection}

\wss{Not required for CAS 741}

\input{../Reflection.text}

\begin{enumerate}
    \item Why is it important to create a development plan prior to starting the
    project?
    \item In your opinion, what are the advantages and disadvantages of using
    CI/CD?
    \item What disagreements did your group have in this deliverable, if any,
    and how did you resolve them?
\end{enumerate}

\newpage{}

\section*{Appendix --- Team Charter}

\wss{borrows from
\href{https://engineering.up.edu/industry_partnerships/files/team-charter.pdf}
{University of Portland Team Charter}}

\subsection*{External Goals}

\wss{What are your team's external goals for this project? These are not the
goals related to the functionality or quality fo the project.  These are the
goals on what the team wishes to achieve with the project.  Potential goals are
to win a prize at the Capstone EXPO, or to have something to talk about in
interviews, or to get an A+, etc.}

\subsection*{Attendance}

\subsubsection*{Expectations}

\wss{What are your team's expectations regarding meeting attendance (being on
time, leaving early, missing meetings, etc.)?}

\subsubsection*{Acceptable Excuse}

\wss{What constitutes an acceptable excuse for missing a meeting or a deadline?
What types of excuses will not be considered acceptable?}

\subsubsection*{In Case of Emergency}

\wss{What process will team members follow if they have an emergency and cannot
attend a team meeting or complete their individual work promised for a team
deliverable?}

\subsection*{Accountability and Teamwork}

\subsubsection*{Quality} 

\wss{What are your team's expectations regarding the quality
of team members' preparation for team meetings and the quality of the
deliverables that members bring to the team?}

\subsubsection*{Attitude}

\wss{What are your team's expectations regarding team members' ideas,
interactions with the team, cooperation, attitudes, and anything else regarding
team member contributions?  Do you want to introduce a code of conduct?  Do you
want a conflict resolution plan?  Can adopt existing codes of conduct.}

\subsubsection*{Stay on Track}

\wss{What methods will be used to keep the team on track? How will your team
ensure that members contribute as expected to the team and that the team
performs as expected? How will your team reward members who do well and manage
members whose performance is below expectations?  What are the consequences for
someone not contributing their fair share?}

\wss{You may wish to use the project management metrics collected for the TA and
instructor for this.}

\wss{You can set target metrics for attendance, commits, etc.  What are the
consequences if someone doesn't hit their targets?  Do they need to bring the
coffee to the next team meeting?  Does the team need to make an appointment with
their TA, or the instructor?  Are there incentives for reaching targets early?}

\subsubsection*{Team Building}

\wss{How will you build team cohesion (fun time, group rituals, etc.)? }

\subsubsection*{Decision Making} 

\wss{How will you make decisions in your group? Consensus?  Vote? How will you
handle disagreements? }

\end{document}